\section{Testy systemu i wyniki analizy}

Po zakończeniu implementacji przeprowadzono testy działania systemu oraz sprawdzono wyniki generowane dla wybranych nagrań martwego ciągu. Celem tego etapu było zweryfikowanie, czy przygotowana aplikacja działa zgodnie z założeniami opisanymi we wcześniejszych rozdziałach. Szczególną uwagę zwrócono na pełny przebieg analizy, czyli przesłanie pliku wideo, uruchomienie przetwarzania po stronie serwera, wykrycie powtórzeń oraz prezentację wyników użytkownikowi.

Testy miały charakter funkcjonalny i jakościowy. Nie przygotowywano dużego zbioru nagrań testowych ani ręcznie oznaczonych danych referencyjnych, dlatego w rozdziale nie przedstawiono statystycznej oceny skuteczności modelu. Takie badanie wymagałoby osobnego zbioru danych oraz porównania wyników systemu z oceną ekspercką lub pomiarem wykonanym inną metodą. W ramach pracy sprawdzono natomiast, czy zaimplementowany system poprawnie obsługuje podstawowy scenariusz użycia oraz jak reaguje na wybrane sytuacje utrudniające analizę obrazu.

W dalszej części rozdziału opisano zakres przeprowadzonych testów, środowisko testowe oraz charakter nagrań wykorzystanych do weryfikacji działania aplikacji. Następnie przedstawiono wyniki analizy nagrań oraz ograniczenia, które zaobserwowano podczas testowania rozwiązania.

\subsection{Cel i zakres testów}

Celem testów było sprawdzenie działania systemu z perspektywy użytkownika oraz potwierdzenie, że poszczególne elementy aplikacji współpracują ze sobą w oczekiwany sposób zgodnie z wymaganiami projektowymi. Skupiono się na testach analizy nagrania w zależności od różnych przypadków w danych wejściowych jak i na sprawdzeniu jak zachowuj się system w momencie napotkania błędów lub złego użytkowania.

Zakres testów obejmował działanie aplikacji sieciowej oraz modułu analizy wideo. W części dotyczącej warstwy aplikacji sieciowej przetestowano możliwość przesłania pliku przez interfejs użytkownika, zapisania informacji o analizie, uruchomienie procesu przetwarzania oraz odczytania wyników analizy po jej zakończeniu. W części dotyczącej analizy skupiono się na pracy reguł badających technikę wykonania ćwiczenia i potencjalne błędy,

Weryfikacja została przeprowadzona na kilku nagraniach przygotowanych na potrzeby pracy. Obejmowały one zarówno poprawne wykonanie ćwiczenia, jak i przykłady problematyczne, takie jak zbyt krótki zakres ruchu, niekorzystne ustawienie kamery, zaokrąglenie pleców oraz ubiór utrudniający rozpoznanie punktów kluczowych sylwetki. Materiał testowy nie miał na celu pokrycia wszystkich możliwych błędów technicznych, lecz pozwolił sprawdzić najważniejsze elementy działania systemu.

Przy interpretacji wyników uwzględniono fakt, że system opiera się na gotowym modelu estymacji pozy człowieka. Oznacza to, że dalsze obliczenia zależą od jakości punktów kluczowych zwróconych przez model. W praktyce istotne znaczenie mają więc takie czynniki jak ustawienie kamery, widoczność całej sylwetki, jakość nagrania oraz stabilność detekcji w kolejnych klatkach filmu.

\subsection{Środowisko testowe i materiał testowy}

Testy systemu przeprowadzono lokalnie na komputerze Lenovo ThinkPad P14s. Moduł analizy wideo uruchamiano w środowisku Python skonfigurowanym w programie PyCharm. Takie środowisko odpowiadało warunkom wykorzystywanym podczas implementacji i pozwalało na bieżąco sprawdzać działanie aplikacji serwerowej oraz modułu odpowiedzialnego za analizę nagrania.

Do estymacji pozy człowieka wykorzystano wariant modelu MediaPipe Pose Landmarker oparty na pliku \texttt{pose\_landmarker\_heavy.task}. Jest to najwolniejszy z dostępnych wariantów, ale jednocześnie zapewnia większą dokładność detekcji punktów kluczowych niż lżejsze modele. W kontekście pracy ważniejsza była jakość uzyskanych danych niż szybkość przetwarzania, dlatego podczas testów wybrano właśnie ten wariant.

Materiał testowy stanowiły nagrania martwego ciągu zarejestrowane telefonem iPhone 14 Pro w rozdzielczości Full HD przy 30 klatkach na sekundę. Część nagrań wykonano zgodnie z założeniami systemu, czyli z ustawieniem kamery z boku, widoczną całą sylwetką oraz względnie stabilnym kadrem. Takie materiały wykorzystano do sprawdzenia podstawowego przebiegu analizy, w tym wykrywania powtórzeń i obliczania parametrów ruchu.

Najważniejsze informacje dotyczące środowiska testowego, użytych technologii oraz parametrów materiału wideo zestawiono w tabeli~\ref{tab:srodowisko-testowe}.

\begin{table}[H]
\centering
\caption{Środowisko testowe i parametry materiału wideo}
\label{tab:srodowisko-testowe}
\begin{tabular}{|
>{\raggedright\arraybackslash}p{4.5cm}|
>{\raggedright\arraybackslash}p{8cm}|}
\hline
\textbf{Element} & \textbf{Opis} \\ \hline
Komputer testowy & Lenovo ThinkPad P14s \\ \hline
Środowisko uruchomieniowe & Lokalny runtime Python skonfigurowany w środowisku PyCharm \\ \hline
Wersja języka Python & Python 3.13.5 \\ \hline
Framework aplikacji serwerowej & Django 6.0.4 \\ \hline
Biblioteka aplikacji klienckiej & React 19.2.6 \\ \hline
Model estymacji pozy & MediaPipe Pose Landmarker, wariant \texttt{pose\_landmarker\_heavy.task} \\ \hline
Urządzenie rejestrujące nagrania & iPhone 14 Pro \\ \hline
Parametry nagrań & Rozdzielczość Full HD, 30 klatek na sekundę \\ \hline
Charakter testów & Testy funkcjonalne i jakościowa ocena wyników analizy \\ \hline
\end{tabular}
\end{table}

Wykorzystano również nagrania zawierające sytuacje utrudniające analizę. Obejmowały one między innymi zbyt krótki zakres ruchu, zaokrąglenie pleców, niekorzystne ustawienie kamery oraz ubiór utrudniający rozpoznanie wybranych punktów kluczowych sylwetki. Przypadki te nie tworzyły osobnego, dużego zbioru testowego. Służyły raczej do sprawdzenia, jak system zachowuje się w warunkach odbiegających od założeń przyjętych podczas projektowania rozwiązania.

Podczas oceny wyników zwracano uwagę przede wszystkim na to, czy nagranie pozwala na wykrycie sylwetki w kolejnych fazach ruchu oraz czy uzyskane punkty kluczowe umożliwiają obliczenie kątów wykorzystywanych w dalszej analizie. Nagrania, w których sylwetka była częściowo zasłonięta, źle ustawiona względem kamery lub niestabilnie wykrywana, traktowano jako materiał pokazujący ograniczenia systemu, a nie jako podstawę do jednoznacznej oceny techniki ćwiczącego.

\subsection{Testy działania aplikacji}

Pierwsza część testów dotyczyła podstawowego scenariusza działania aplikacji. Sprawdzono proces, w którym użytkownik przesyła nagranie martwego ciągu, system uruchamia analizę po stronie serwera, a po zakończeniu przetwarzania wyświetla wynik w aplikacji klienckiej. Test obejmował więc pełną ścieżkę obsługi nagrania, od wyboru pliku do prezentacji rezultatów.

Po wybraniu pliku wideo w interfejsie użytkownika nagranie było przesyłane do aplikacji serwerowej. Serwer zapisywał plik oraz tworzył rekord analizy zawierający informacje potrzebne do dalszego przetwarzania. Następnie uruchamiany był moduł analizy martwego ciągu, który przetwarzał nagranie i zwracał dane obejmujące wykryte powtórzenia, podsumowanie parametrów ruchu oraz ocenę techniki.

Po zakończeniu analizy wynik był zapisywany po stronie serwera i udostępniany aplikacji klienckiej. W interfejsie użytkownika sprawdzono, czy widoczna jest liczba wykrytych powtórzeń, ocena poszczególnych ruchów oraz szczegóły zawierające wybrane parametry biomechaniczne. Pozwoliło to potwierdzić, że dane wygenerowane przez moduł analizy są poprawnie przekazywane między backendem i frontendem.

Przykład widoku, w którym prezentowany jest wynik zakończonej analizy, przedstawiono wcześniej na rysunku~\ref{fig:widok-wynikow-analizy}.

W testowanym scenariuszu aplikacja wykonała pełny proces analizy nagrania bez konieczności ręcznej ingerencji w dane pośrednie. Użytkownik przekazywał jedynie plik wideo, a pozostałe etapy były realizowane po stronie systemu. Otrzymany wynik był dostępny w aplikacji klienckiej w formie przewidzianej w projekcie interfejsu.

Test podstawowej ścieżki działania potwierdził poprawną współpracę głównych elementów systemu. Aplikacja kliencka umożliwiała przesłanie nagrania, aplikacja serwerowa uruchamiała analizę, a wynik zwracany przez moduł przetwarzania wideo był prezentowany użytkownikowi. Na tej podstawie możliwe było przejście do testów dotyczących jakości samej analizy nagrań.

\subsection{Testy obsługi błędnych danych wejściowych}

Oprócz podstawowej ścieżki działania sprawdzono również zachowanie aplikacji w sytuacji, gdy użytkownik wybiera plik w nieobsługiwanym formacie. Taki przypadek jest istotny, ponieważ moduł analizy został przygotowany do przetwarzania nagrań wideo, a przekazanie innego typu pliku nie powinno prowadzić do uruchomienia analizy.

Test polegał na wybraniu w aplikacji klienckiej pliku, który nie był plikiem wideo. W takiej sytuacji interfejs użytkownika nie przesyłał pliku do dalszego przetwarzania i wyświetlał komunikat informujący o błędzie. Treść komunikatu miała postać: \texttt{Błąd: Wybrany plik nie jest plikiem video.}

Sprawdzenie to potwierdziło, że aplikacja obsługuje podstawowy błąd wynikający z nieprawidłowego użycia formularza przesyłania nagrania. Użytkownik otrzymuje informację o przyczynie problemu jeszcze przed uruchomieniem analizy, co ogranicza ryzyko wykonania niepotrzebnego przetwarzania po stronie serwera.

Test ten nie obejmował pełnej walidacji wszystkich możliwych przypadków błędnych danych wejściowych. Jego celem było sprawdzenie podstawowej reakcji aplikacji na sytuację, w której użytkownik próbuje przekazać do systemu plik niezgodny z przeznaczeniem formularza.

\subsection{Wyniki analizy przykładowego nagrania}

Do dokładniejszej analizy wybrano nagranie przedstawiające serię ośmiu powtórzeń martwego ciągu. Plik miał rozmiar około 33 MB, a czas trwania nagrania wynosił 32 sekundy. Materiał został wykonany z boku osoby ćwiczącej, przy stabilnym ustawieniu kamery oraz widocznej całej sylwetce. Ubiór nie utrudniał znacząco analizy, ponieważ kolana i ułożenie tułowia pozostawały widoczne w kolejnych fazach ruchu. W tle nagrania znajdowało się lustro oraz inne osoby ćwiczące, jednak nie zakłóciło to wyboru analizowanej sylwetki.

Przed uruchomieniem analizy nagranie oceniono wizualnie. Przyjęto, że powtórzenia pierwsze, drugie, trzecie, czwarte, siódme i ósme zostały wykonane poprawnie. Różniły się one głównie głębokością pozycji startowej. Pierwsze dwa powtórzenia rozpoczynały się z wyższej pozycji, trzecie oraz ósme z pozycji zbliżonej do środka serii, natomiast czwarte i siódme z pozycji głębszej. Piąte powtórzenie zostało wykonane w zbyt krótkim zakresie ruchu. Szóste powtórzenie miało pełny zakres ruchu, ale podczas jego wykonania widoczne było niekorzystne ustawienie pleców.

Zestawienie oceny wizualnej z wynikiem zwróconym przez system przedstawiono w tabeli~\ref{tab:porownanie-oceny-powtorzen}.

\begin{table}[H]
\centering
\caption{Porównanie oceny wizualnej z wynikiem systemu dla przykładowego nagrania}
\label{tab:porownanie-oceny-powtorzen}
\footnotesize
\renewcommand{\arraystretch}{1.08}
\setlength{\tabcolsep}{2.5pt}
\begin{tabular}{|
>{\centering\arraybackslash}p{1.05cm}|
>{\raggedright\arraybackslash}p{4.25cm}|
>{\raggedright\arraybackslash}p{3.95cm}|
>{\centering\arraybackslash}p{1.65cm}|
>{\raggedright\arraybackslash}p{3.85cm}|}
\hline
\textbf{Powt.} & \textbf{Ocena wizualna} & \textbf{Wynik systemu} & \textbf{Zgodność} & \textbf{Komentarz} \\ \hline
1 & Poprawne, pełny zakres ruchu, dobry wyprost bez widocznego przeprostu, wyższa pozycja startowa & Poprawne, wyższa pozycja startowa & Tak & System poprawnie rozpoznał charakter pozycji startowej i nie wskazał błędów techniki. \\ \hline
2 & Poprawne, pełny zakres ruchu, dobry wyprost bez widocznego przeprostu, wyższa pozycja startowa & Poprawne, wyższa pozycja startowa & Tak & Wynik systemu odpowiadał ocenie wizualnej. \\ \hline
3 & Poprawne, pełny zakres ruchu, jednostajny przebieg ruchu, pozycja startowa średnia & Poprawne, pozycja startowa zbliżona do mediany serii & Tak & System wskazał pozycję startową zgodną z obserwacją nagrania. \\ \hline
4 & Poprawne, pełny zakres ruchu, dobry wyprost bez widocznego przeprostu, głębsza pozycja startowa & Poprawne, głębsza pozycja startowa & Tak & System poprawnie odróżnił głębszą pozycję startową od wcześniejszych powtórzeń. \\ \hline
5 & Niepoprawne, zbyt krótki zakres ruchu, brak pełnego wyprostu w pozycji końcowej & Niepoprawne, obniżony zakres ruchu oraz niepełny wyprost biodra i kolana & Tak & System wykrył główny błąd widoczny w nagraniu. \\ \hline
6 & Wymaga uwagi, pełny zakres ruchu, widoczne zaokrąglenie pleców & Wymaga uwagi, bark niżej niż biodro w pozycji startowej & Częściowo & System wskazał wątpliwe ustawienie tułowia, nie samo przegięcie pleców. \\ \hline
7 & Poprawne, pełny zakres ruchu, dobry wyprost bez widocznego przeprostu, głębsza pozycja startowa & Poprawne, głębsza pozycja startowa & Tak & Wynik był zgodny z oceną wizualną. \\ \hline
8 & Poprawne, pełny zakres ruchu, jednostajny przebieg ruchu, pozycja startowa średnia & Poprawne, pozycja startowa zbliżona do mediany serii & Tak & System poprawnie ocenił powtórzenie i jego pozycję startową. \\ \hline
\end{tabular}
\end{table}

Przykładową klatkę z analizowanego nagrania, wraz z naniesionymi punktami kluczowymi sylwetki oraz połączeniami wykorzystywanymi podczas analizy, przedstawiono na rysunku~\ref{fig:klatka-analizy-przykladowe-nagranie}.

\begin{figure}[H]
\centering
\includegraphics[width=0.55\textwidth]{figures/obraz_punkty_kluczowe.png}
\caption{Przykładowa klatka z analizowanego nagrania z naniesionymi punktami kluczowymi sylwetki wykrytymi przez model estymacji pozy. Źródło: opracowanie własne}
\label{fig:klatka-analizy-przykladowe-nagranie}
\end{figure}
System wykrył w nagraniu osiem powtórzeń, co odpowiadało rzeczywistej liczbie ruchów w serii. Powtórzenia od pierwszego do czwartego zostały oznaczone jako poprawne. Wynik ten był zgodny z oceną wizualną. W tych powtórzeniach zakres ruchu sztangi był zbliżony do mediany serii, a w górnej pozycji odnotowano wyprost biodra i kolana mieszczący się w przyjętych kryteriach. System poprawnie rozróżnił także charakter pozycji startowej. Pierwsze dwa powtórzenia zostały opisane jako rozpoczynające się z wyższej pozycji, trzecie jako zbliżone do mediany serii, a czwarte jako głębsze.

Piąte powtórzenie zostało sklasyfikowane jako niepoprawne. Był to wynik zgodny z obserwacją nagrania, ponieważ ruch nie został wykonany w pełnym zakresie. System wskazał trzy przyczyny: niepełny wyprost biodra, niepełny wyprost kolana oraz obniżony zakres ruchu sztangi. Wartość zakresu ruchu sztangi dla tego powtórzenia wyniosła 0{,}110, podczas gdy mediana serii wynosiła 0{,}161. Różnica ta dobrze oddaje charakter błędu widocznego na nagraniu. Dodatkowo w górnej pozycji kąt biodra wyniósł 80{,}6°, a kąt kolana 120{,}9°, co potwierdza, że ruch zakończył się przed osiągnięciem pełnej pozycji końcowej.

Szóste powtórzenie jest najważniejsze z punktu widzenia oceny ograniczeń systemu. W ocenie wizualnej było ono problematyczne ze względu na zaokrąglenie pleców, mimo że zakres ruchu był pełny. System nie sklasyfikował tego ruchu jako niepoprawnego, ponieważ wartości zakresu ruchu oraz wyprostu w górnej pozycji mieściły się w przyjętych kryteriach. Jednocześnie powtórzenie zostało oznaczone w interfejsie jako wymagające uwagi. W szczegółowym komunikacie wskazano, że w pozycji startowej bark znajdował się niżej niż biodro, co może sugerować niekorzystne ustawienie tułowia lub podejrzenie zaokrąglenia pleców. W tym przypadku system poprawnie zasygnalizował problem z pozycją pleców, ale potraktował go jako ostrzeżenie, a nie jako błąd jednoznacznie unieważniający całe powtórzenie.

Przykładową klatkę z szóstego powtórzenia, w którym system wskazał problem z relacją barku i biodra w pozycji startowej, przedstawiono na rysunku~\ref{fig:klatka-powtorzenie-szoste}.

\begin{figure}[H]
\centering
\includegraphics[width=0.55\textwidth]{figures/obraz_koci_grzebiet.png}
\caption{Klatka z szóstego powtórzenia, w którym system wskazał niekorzystne ustawienie tułowia w pozycji startowej. Źródło: opracowanie własne}
\label{fig:klatka-powtorzenie-szoste}
\end{figure}

Na rysunku widoczne jest obniżenie barków względem bioder oraz zaokrąglenie kręgosłupa w odcinku lędźwiowym, tak zwany \enquote{Koci Grzbiet}, Na tej podstawie system wygenerował ostrzeżenie co do pozycji barku względem biodra.

Powtórzenia siódme i ósme zostały ocenione jako poprawne. Dla siódmego powtórzenia system wskazał głębszą pozycję startową, co było zgodne z obserwacją nagrania. Ósme powtórzenie zostało opisane jako rozpoczynające się z pozycji zbliżonej do mediany serii. Oba powtórzenia miały zakres ruchu zbliżony do pozostałych poprawnych ruchów oraz spełniły kryteria wyprostu w końcowej fazie.

Oprócz oceny pojedynczych powtórzeń system przygotował także ocenę spójności serii. Ogólny status tej oceny został oznaczony jako ostrzeżenie. Zakres ruchu sztangi między poprawnymi powtórzeniami był stabilny, a różnica wyniosła 0{,}022. Większą zmienność odnotowano dla czasu trwania powtórzeń oraz kąta kolana w górnej pozycji. Różnica czasu trwania poprawnych powtórzeń wyniosła 2{,}53 s, natomiast różnica kąta kolana w pozycji końcowej wyniosła 12{,}0°. Ostrzeżenie to jest uzasadnione, ponieważ szóste powtórzenie trwało wyraźnie dłużej od pozostałych, a część ruchów różniła się sposobem dojścia do pozycji końcowej.

Analiza przykładowego nagrania potwierdziła, że system poprawnie wykrywa liczbę powtórzeń oraz potrafi wskazać ruch o ograniczonym zakresie. Wyniki dla pozycji startowej były w dużej mierze zgodne z oceną wizualną. W przypadku powtórzenia z zaokrągleniem pleców system nie oznaczył ruchu jako niepoprawnego, ale poprawnie wskazał go jako wymagający uwagi. Jest to zgodne z przyjętym sposobem działania aplikacji, w którym część problemów technicznych ma charakter ostrzeżenia, a nie automatycznego unieważnienia powtórzenia.

\subsection{Analiza nagrania wykonanego pod niekorzystnym kątem}

Pierwszy przypadek problematyczny dotyczył nagrania wykonanego pod kątem względem osoby ćwiczącej. Materiał przedstawiał sześć powtórzeń martwego ciągu. Cztery pierwsze ruchy oceniono wizualnie jako poprawne, natomiast dwa ostatnie zawierały widoczne zaokrąglenie pleców. Sylwetka była widoczna w całości, jednak nagranie nie spełniało w pełni założenia widoku bocznego. Kamera była ustawiona skośnie i nieznacznie zmieniała swoje położenie. W tle widoczne były także lustra oraz inne osoby ćwiczące.

System wykrył pięć powtórzeń. Różnica względem rzeczywistej liczby ruchów dotyczyła początku nagrania. Pierwsze powtórzenie zostało rozpoczęte bezpośrednio po uruchomieniu nagrywania i w logach analizy było traktowane jako fragment bez przypisanego powtórzenia. Z tego powodu nie znalazło się w końcowym podsumowaniu. Prawdopodobnie wpływ miało tutaj zarówno zbyt szybkie rozpoczęcie ruchu po starcie filmu, jak i niekorzystne ustawienie kamery. Przy skośnym ujęciu zmiana położenia punktów kluczowych oraz przybliżonej pozycji sztangi może być mniej jednoznaczna, co utrudnia wykrycie pełnego cyklu ruchu. Pokazuje to, że system najlepiej działa wtedy, gdy nagranie zawiera krótki fragment poprzedzający pierwsze powtórzenie oraz spełnia założenie bocznego ustawienia kamery.

Porównanie oceny wizualnej z wynikiem systemu dla tego nagrania przedstawiono w tabeli~\ref{tab:porownanie-nagranie-pod-katem}.

\begin{table}[H]
\centering
\caption{Porównanie oceny wizualnej z wynikiem systemu dla nagrania wykonanego pod kątem}
\label{tab:porownanie-nagranie-pod-katem}
\footnotesize
\renewcommand{\arraystretch}{1.08}
\setlength{\tabcolsep}{2.5pt}
\begin{tabular}{|
>{\centering\arraybackslash}p{1.0cm}|
>{\raggedright\arraybackslash}p{3.7cm}|
>{\raggedright\arraybackslash}p{3.8cm}|
>{\centering\arraybackslash}p{1.6cm}|
>{\raggedright\arraybackslash}p{4.0cm}|}
\hline
\textbf{Powt.} & \textbf{Ocena wizualna} & \textbf{Wynik systemu} & \textbf{Zgodność} & \textbf{Komentarz} \\ \hline
1 & Poprawne, pełny ruch & Niewykryte w podsumowaniu & Nie & Boczna perspektywa i szybkie zaczęcie ruchu wpłyneło na brak wykrycia. \\ \hline
2 & Poprawne, pełny zakres ruchu, prawidłowy wyprost & Poprawne, pozycja startowa zbliżona do mediany serii & Tak & System poprawnie ocenił zakres ruchu oraz pozycję końcową. \\ \hline
3 & Poprawne, pełny zakres ruchu, prawidłowy wyprost & Poprawne, pozycja startowa zbliżona do mediany serii & Tak & Wynik systemu odpowiadał ocenie wizualnej. \\ \hline
4 & Poprawne, pełny zakres ruchu, prawidłowy wyprost & Częściowe, niepełny wyprost kolana & Częściowo & System wskazał problem, który nie był jednoznacznie widoczny podczas oceny wizualnej. \\ \hline
5 & Wymaga uwagi, widoczne zaokrąglenie pleców & Poprawne, wyższa pozycja startowa & Nie & System nie wykrył problemu z ustawieniem pleców, mimo że ruch był problematyczny wizualnie. \\ \hline
6 & Wymaga uwagi, widoczne zaokrąglenie pleców & Poprawne, pozycja startowa zbliżona do mediany serii & Nie & System poprawnie ocenił zakres ruchu, ale nie wskazał ostrzeżenia dotyczącego tułowia. \\ \hline
\end{tabular}
\end{table}

Dla wykrytych powtórzeń system poprawnie rozpoznał zakres ruchu sztangi. Wartości tego parametru były zbliżone do mediany serii, a ocena spójności zakresu ruchu została oznaczona jako poprawna. Również wyprost w górnej pozycji został w większości przypadków oceniony pozytywnie. Z punktu widzenia podstawowej detekcji ruchu system zachował więc stabilność mimo niekorzystnego ustawienia kamery.

Większe różnice pojawiły się przy interpretacji kątów i ustawienia tułowia. Trzecie wykryte powtórzenie zostało oznaczone jako częściowe ze względu na niepełny wyprost kolana. W ocenie wizualnej powtórzenie to nie wyróżniało się jednak wyraźnym błędem technicznym. Można więc przyjąć, że wynik ten był związany z perspektywą nagrania. Przy ustawieniu kamery pod kątem położenie kolana, biodra i barku jest rzutowane inaczej niż w ujęciu bocznym, na którym opiera się przyjęty sposób analizy.

Najważniejsza rozbieżność dotyczyła końcowych powtórzeń. Dwa ostatnie ruchy były wizualnie problematyczne ze względu na zaokrąglenie pleców. System nie oznaczył ich jednak jako wymagających uwagi w zakresie ustawienia tułowia. W komunikatach szczegółowych relacja barku i biodra została oceniona jako poprawna. Oznacza to, że w tym przypadku model oraz reguły oceny nie uchwyciły błędu widocznego dla obserwatora na nagraniu.

Klatkę z jednego z końcowych powtórzeń, w którym podczas oceny wizualnej zauważono zaokrąglenie pleców, przedstawiono na rysunku~\ref{fig:nagranie-pod-katem-plecy}.

\begin{figure}[H]
\centering
\includegraphics[width=0.55\textwidth]{figures/obraz_boczna.png}
\caption{Klatka z nagrania wykonanego pod niekorzystnym kątem, przedstawiająca powtórzenie z widocznym problemem ustawienia pleców. Źródło: opracowanie własne}
\label{fig:nagranie-pod-katem-plecy}
\end{figure}

Przedstawiony przykład pokazuje, że przy skośnym ustawieniu kamery błąd widoczny podczas oceny wizualnej nie musi zostać jednoznacznie odzwierciedlony w relacjach pomiędzy punktami kluczowymi wykorzystywanymi przez system.

Wynik ten nie oznacza błędnego działania całej aplikacji, lecz pokazuje ograniczenie przyjętej metody. System analizuje obraz dwuwymiarowy i zakłada, że nagranie przedstawia osobę ćwiczącą z boku. Jeżeli kamera jest ustawiona skośnie, obliczane kąty oraz relacje pomiędzy punktami kluczowymi mogą nie odpowiadać rzeczywistemu ustawieniu ciała. W takim przypadku wynik dotyczący zakresu ruchu może nadal być użyteczny, ale ocena ustawienia pleców i tułowia wymaga większej ostrożności.

Analiza tego nagrania potwierdziła znaczenie wymagań dotyczących sposobu przygotowania materiału wejściowego. Dla poprawnej pracy systemu istotne jest nie tylko to, aby cała sylwetka była widoczna, ale również to, aby kamera była ustawiona możliwie prostopadle do kierunku ruchu. Ważne jest także rozpoczęcie nagrania chwilę przed pierwszym powtórzeniem, tak aby system mógł uchwycić pełny przebieg ruchu od pozycji startowej.

\subsection{Nagranie z ubiorem utrudniającym detekcję kolan}

Kolejny przypadek problematyczny dotyczył nagrania, w którym osoba ćwicząca miała na sobie szerokie, jednolite, czarne spodnie. Materiał przedstawiał pięć poprawnie wykonanych powtórzeń martwego ciągu. Kamera była ustawiona w sposób umożliwiający obserwację całej sylwetki, a ruch został wykonany w pełnym zakresie. Podczas oceny wizualnej nie stwierdzono istotnych błędów techniki w żadnym z powtórzeń.

Mimo poprawnego wykonania ćwiczenia nagranie było trudniejsze dla modelu estymacji pozy. Szerokie spodnie utrudniały jednoznaczne określenie położenia kolan, szczególnie w fazie opuszczania ciężaru. Przykład błędnej lokalizacji punktu kolanowego w pojedynczej klatce nagrania przedstawiono na rysunku~\ref{fig:bledna-detekcja-kolana}.

\begin{figure}[H]
\centering
\includegraphics[width=0.55\textwidth]{figures/obraz_spodnie.png}
\caption{Przykład błędnej detekcji punktu kolanowego w nagraniu z ubiorem utrudniającym rozpoznanie położenia stawu. Źródło: opracowanie własne}
\label{fig:bledna-detekcja-kolana}
\end{figure}

Na przedstawionej klatce punkt kolanowy został wykryty z wyraźnym przesunięciem względem rzeczywistego położenia stawu, co może wpływać na obliczane wartości kąta kolana w dalszej analizie.

System wykrył w nagraniu pięć powtórzeń, co odpowiadało rzeczywistej liczbie ruchów. Wszystkie powtórzenia zostały oznaczone jako poprawne. Zakres ruchu sztangi został oceniony jako prawidłowy, a wartości tego parametru były zbliżone do mediany serii. W każdym powtórzeniu system wskazał również poprawny wyprost biodra oraz kolana w pozycji końcowej. Z punktu widzenia ogólnej klasyfikacji serii wynik był więc zgodny z oceną wizualną.

Porównanie oceny wizualnej z wynikiem zwróconym przez system przedstawiono w tabeli~\ref{tab:porownanie-ubior-kolana}.

\begin{table}[H]
\centering
\caption{Porównanie oceny wizualnej z wynikiem systemu dla nagrania z ubiorem utrudniającym detekcję kolan}
\label{tab:porownanie-ubior-kolana}
\footnotesize
\renewcommand{\arraystretch}{1.08}
\setlength{\tabcolsep}{2.5pt}
\begin{tabular}{|
>{\centering\arraybackslash}p{1.05cm}|
>{\raggedright\arraybackslash}p{4.15cm}|
>{\raggedright\arraybackslash}p{4.05cm}|
>{\centering\arraybackslash}p{1.55cm}|
>{\raggedright\arraybackslash}p{3.85cm}|}
\hline
\textbf{Powt.} & \textbf{Ocena wizualna} & \textbf{Wynik systemu} & \textbf{Zgodność} & \textbf{Komentarz} \\ \hline
1 & Poprawne, pełny zakres ruchu, prawidłowy wyprost & Poprawne, pozycja startowa zbliżona do mediany serii & Tak & System poprawnie ocenił ruch i nie wskazał błędów techniki. \\ \hline
2 & Poprawne, pełny zakres ruchu, prawidłowy wyprost & Poprawne, pozycja startowa zbliżona do mediany serii & Tak & Wynik systemu odpowiadał ocenie wizualnej. \\ \hline
3 & Poprawne, pełny zakres ruchu, prawidłowy wyprost & Poprawne, pozycja startowa zbliżona do mediany serii & Tak & System poprawnie wykrył powtórzenie mimo trudniejszej widoczności kolan. \\ \hline
4 & Poprawne, pełny zakres ruchu, prawidłowy wyprost & Poprawne, pozycja startowa zbliżona do mediany serii & Tak & Ogólna klasyfikacja powtórzenia była zgodna z oceną wizualną. \\ \hline
5 & Poprawne, pełny zakres ruchu, prawidłowy wyprost & Poprawne, wyższa pozycja startowa & Częściowo & System poprawnie ocenił powtórzenie jako poprawne, ale odmiennie zinterpretował pozycję startową, prawdopodobnie przez mniej stabilne wykrycie kolana. \\ \hline
\end{tabular}
\end{table}

Wpływ problematycznego ubioru był jednak widoczny w danych szczegółowych. Największe różnice pojawiły się w parametrach zależnych od położenia kolana. W piątym powtórzeniu system opisał pozycję startową jako wyższą, wskazując kąt kolana równy 133{,}7°. Pozostałe powtórzenia zostały opisane jako rozpoczynające się z pozycji zbliżonej do mediany serii. Tak duża różnica wartości kąta kolana może wynikać nie tylko z rzeczywistej zmiany pozycji, ale również z mniej stabilnego wykrycia punktu kolanowego.

Podobny problem widoczny był w ocenie spójności serii. Ogólny status spójności został oznaczony jako ostrzeżenie. Zakres ruchu sztangi oraz wyprost biodra pozostały stabilne, natomiast dla kąta kolana w górnej pozycji system wskazał różnicę 12{,}9°. Jest to zgodne z obserwacją, że w tym nagraniu największą niepewność wprowadzały punkty związane z kończynami dolnymi. Warto zauważyć, że pojedyncze błędne klatki nie spowodowały błędnej klasyfikacji całych powtórzeń, ale wpłynęły na ocenę parametrów szczegółowych i spójności serii.

Analiza tego przypadku pokazuje, że system może poprawnie rozpoznać liczbę powtórzeń i ogólny przebieg ruchu nawet wtedy, gdy wybrane punkty kluczowe są w części klatek wykrywane niestabilnie. Jednocześnie wyniki dotyczące kątów kolana należy w takich warunkach interpretować ostrożnie. Ubiór zasłaniający lub zlewający się z tłem może obniżać wiarygodność pomiarów, szczególnie wtedy, gdy punkt stawu jest dodatkowo zasłonięty przez sztangę albo inną część ciała.

\subsection{Podsumowanie wyników testów}

Przeprowadzone testy potwierdziły, że zaimplementowany system realizuje główny scenariusz działania przyjęty w pracy. Aplikacja umożliwia przesłanie nagrania wideo, uruchamia analizę po stronie serwera, zapisuje uzyskane wyniki oraz prezentuje je użytkownikowi w interfejsie aplikacji klienckiej. Sprawdzono również podstawową obsługę błędnych danych wejściowych, w której aplikacja informuje użytkownika o wyborze pliku niebędącego nagraniem wideo.

Najlepsze wyniki uzyskano dla nagrania wykonanego zgodnie z założeniami systemu. Przy ustawieniu kamery z boku, widocznej całej sylwetce oraz ubiorze nieutrudniającym detekcji punktów kluczowych system poprawnie wykrył liczbę powtórzeń i w większości przypadków zwrócił ocenę zgodną z oceną wizualną. W szczególności prawidłowo rozpoznano powtórzenie wykonane w zbyt krótkim zakresie ruchu. System wskazał w nim obniżony zakres ruchu sztangi oraz brak pełnego wyprostu biodra i kolana.

Wyniki pokazały również, że system może sygnalizować problemy z ustawieniem tułowia. W przypadku powtórzenia z widocznym zaokrągleniem pleców aplikacja oznaczyła ruch jako wymagający uwagi i wskazała niekorzystną relację barku względem biodra w pozycji startowej. Oznacza to, że przy spełnionych warunkach nagrania zaimplementowane reguły mogą wspierać użytkownika w wykrywaniu wybranych błędów technicznych.

Przypadki problematyczne potwierdziły jednak, że wiarygodność analizy jest silnie zależna od jakości materiału wejściowego. Nagranie wykonane pod niekorzystnym kątem utrudniło poprawne wykrycie pełnej liczby powtórzeń oraz ocenę ustawienia pleców. Z kolei szerokie, jednolite spodnie powodowały niestabilną lokalizację punktów kolanowych w pojedynczych klatkach. Wpływało to przede wszystkim na wartości kątów zależnych od położenia kolana oraz na ocenę spójności serii.

Na podstawie testów można stwierdzić, że system spełnia założenia pracy jako prototyp aplikacji do analizy martwego ciągu na podstawie nagrania wideo. Rozwiązanie poprawnie obsługuje główny przepływ danych i potrafi wskazać wybrane błędy techniczne w warunkach zgodnych z przyjętymi wymaganiami. Jednocześnie wyniki analizy powinny być interpretowane z uwzględnieniem ograniczeń metody, szczególnie w przypadku nagrań wykonanych pod złym kątem, z częściowo zasłoniętą sylwetką lub z ubiorem utrudniającym detekcję punktów kluczowych.

System nie zastępuje oceny trenera ani profesjonalnej analizy biomechanicznej. Może jednak pełnić funkcję narzędzia wspierającego wstępną ocenę techniki, szczególnie wtedy, gdy użytkownik przygotuje nagranie zgodnie z wymaganiami opisanymi w pracy.